\section{Introduction}\label{sec:introduction}

A Fano manifold is a smooth connected complex projective variety $X$
whose anticanonical divisor $-K_X$ is ample. Its Picard number
$\rho_X$ is the rank of its N\'eron--Severi group, or equivalently the
dimension of the space of numerical divisor classes. If $X$ has positive
dimension, its \emph{pseudoindex} is
\[
 \iota_X=\min\{-K_X\cdot C:C\subset X
                  \text{ is an irreducible rational curve}\}.
\]
The existence of rational curves on a Fano manifold makes this minimum
well defined. The pseudoindex records the smallest anticanonical curve
degree; it need not equal the largest integer dividing $-K_X$ in
$\operatorname{Pic}(X)$, which is the Fano index.

The generalized Mukai conjecture asks for a sharp bound on the Picard
number in terms of dimension and pseudoindex, together with a
classification of the extremal manifolds. We prove the conjectured bound
and classification in arbitrary dimension.

\begin{theorem}\label{thm:main}
Let $X$ be a smooth connected complex projective Fano manifold of
positive dimension $n$, Picard number $\rho_X$ and pseudoindex $\iota_X$.
Then
\[
 \rho_X(\iota_X-1)\le n.
\]
Equality holds if and only if
\[
 X\cong(\mathbb P^{\iota_X-1})^{\rho_X}.
\]
\end{theorem}

The index version originates in Mukai's work~\cite{Mukai1988}.
Wi\'{s}niewski introduced the pseudoindex and proved that
$2\iota_X>n+2$ forces $\rho_X=1$~\cite[Theorem~A]{Wisniewski1990}.
Bonavero, Casagrande, Debarre and Druel formulate the pseudoindex
strengthening and establish it for Fano manifolds of dimension at most
four, and for smooth toric Fano manifolds when $n\le7$ or
$3\iota_X\ge n+3$~\cite{BCDD2003}. Andreatta, Chierici and Occhetta
extend the low-dimensional result to dimension five and prove
high-pseudoindex cases subject to contraction hypotheses
\cite{AndreattaChiericiOcchetta2004}.

For varieties with additional symmetry,
Casagrande settles the bound and equality classification for
$\mathbb Q$-factorial Gorenstein toric Fano varieties in every
dimension~\cite[Theorem~1(ii)]{Casagrande2006}, and Pasquier treats
$\mathbb Q$-factorial Gorenstein horospherical Fano
varieties~\cite[Theorem~1]{Pasquier2010}. More recently, Gagliardi,
Hofscheier and Pearson prove the generalized conjecture for
$\mathbb Q$-factorial spherical Fano varieties
\cite{GagliardiHofscheierPearson2025}.

Large pseudoindex gives a different route. Novelli and Occhetta
prove both conclusions for all smooth Fano manifolds with
$\iota_X\ge(n+3)/3$~\cite{NovelliOcchetta2010}, and Novelli later
enlarges this range to $\iota_X>n/3$~\cite[Theorem~5.1]{Novelli2012}.
At the boundary $n=3\iota_X$, Novelli also proves the conjecture
when $X$ admits an extremal contraction of fiber type
\cite{Novelli2026}. Theorem~\ref{thm:main}
imposes no restriction on the dimension, Picard number or geometry
of the smooth Fano manifold.

The ordered-chain method of Bonavero, Casagrande, Debarre and Druel
derives the inequality once $\rho_X$ numerically independent
proper irreducible components of the space of rational curves, joined
by a nonempty chain, are available~\cite[Corollary~5.3]{BCDD2003}.
Our argument first detects
independent curve degrees in quantum correspondences. Proper covering
families enter only at equality, when all the selected degrees are minimal.

\subsection*{The argument}

Write $r=\rho_X$ and $\iota=\iota_X$. If $\iota=1$, the inequality
is strict because $n>0$, so assume $\iota\ge2$.
Choose divisor classes $D_1,\ldots,D_r$ forming a numerical basis.
Small quantum multiplication by these divisors gives commuting
matrices $A_1(q),\ldots,A_r(q)$ on the even cohomology of $X$,
with coefficients depending on $r$ nonzero formal variables $q_j$.
Each matrix consists of classical cup product and a finite sum of
two-point curve correspondences. A correspondence of anticanonical
degree $d$ lowers cohomological codimension by $d-1$.
Consequently, a nonzero product in $r$ independent curve degrees would
give
\[
 r(\iota-1)\le\sum_{j=1}^r(d_j-1)\le n.
\]
The main task is to produce that product without assuming that quantum
cohomology is semisimple.

We first prove a lower bound for a point descendant, an intersection
number on the space of one-pointed rational stable maps involving the
$(d-2)$nd power of the marked cotangent class. Here $d$ is the
anticanonical degree. For any degree represented by a free rational map
(a map whose pulled-back tangent bundle is globally generated), this
number is at least $(ad)^{-(d-1)}$, where $-aK_X$
is a fixed very ample divisor. To prove positivity, we contract all
components away from the marking and map the evaluation fiber to a
weighted projective stack. The image of its boundary has dimension
too small to contribute. The pushed-forward virtual class is therefore
a nonempty effective cycle, whose degree admits the required uniform
bound. Section~\ref{sec:descendants} gives the construction in families
and keeps track of stack multiplicities.

Two-point descendant solutions of quantum differential equations
provide a conceptual antecedent for the recurrence used next
\cite[\S6]{Givental1996}.
Here the divisor and recursion identities give a finite-shift
recurrence for descendant vectors on an arbitrary smooth Fano manifold.
After factorial normalization it yields an
exponential upper bound. If every joint eigenvalue tuple of the quantum
divisors satisfied a polynomial relation, the same recurrence would
force decay of order $\exp(-c d\log d)$ along a suitable free curve
degree. After the same factorial normalization, the descendant lower
bound permits only exponential decay and rules this out. Thus one joint
tuple has $r$ algebraically independent coordinates
(Section~\ref{sec:spectral}).

Section~\ref{sec:trace} converts this algebraic independence into the
required curve correspondences. An alternating trace of matrix-valued
differentials is invariant even under changes of basis depending on
$q$. Computing it in a common triangular basis detects the independent
eigenvalues; computing it in the fixed cohomology basis detects a
nonzero product in independent curve degrees. This argument also
applies to nonsemisimple commuting matrices.

At equality all these degrees are minimal. Their nonzero product
implies the existence of a chain of actual rational curves; minimality
makes the full irreducible families containing them proper.
The ordered-chain bound of Bonavero, Casagrande, Debarre and Druel
applies to these components, as required by its erratum
\cite{BCDDErratum}. A cyclic reordering makes every family covering,
and Occhetta's characterization of products of projective spaces gives
the asserted rigidity~\cite{Occhetta2006}.
Section~\ref{sec:equality} verifies each of these family hypotheses.
